\section{Representation Learning：Contrastive 与 Predictive Objectives}

有了离散 code 后，第二篇论文比较两种不依赖标签的路线。Contrastive learning 让同一声音的不同 augmentation 靠近、不同声音远离；generative/predictive learning 则要求 latent code 预测未来。两者都希望学到可迁移 representation，但保留的信息不同：前者编码所定义的不变性，后者编码对未来可预测性有用的状态。

\subsection{Contrastive Learning 在拉近什么}

对 anchor $x$、positive augmentation $x^+$ 与一组 negatives $x_j^-$，可写一个 InfoNCE-style objective：
\[
\mathcal{L}_{\mathrm{NCE}}
=-\log
\frac{\exp(\mathrm{sim}(f(x),f(x^+))/\tau)}
{\exp(\mathrm{sim}(f(x),f(x^+))/\tau)
+\sum_j\exp(\mathrm{sim}(f(x),f(x_j^-))/\tau)}.
\]
其中 $f$ 是 encoder，$\mathrm{sim}$ 是相似度，$\tau$ 是 temperature。Augmentation 不是中性操作：若 pitch shift 被当作 positive，表示会更忽略音高；这对 scene classification 可能有利，对 pitch estimation 却可能有害。

\lecturefigure{slide-32-contrastive-architecture.jpg}{Contrastive audio representation：augmentation pairs、encoder 与相似度目标。}{Stanford Online 官方视频 00:30:25--00:32:25。}

\begin{warningbox}{Invariant to what?}
Contrastive objective 不会自动发现“正确语义”，它会学习 augmentation 定义的不变性。Audio 中 pitch、tempo、timbre、noise 与 time shift 是否应被忽略，取决于下游任务；错误 augmentation 会主动删除关键信息。
\end{warningbox}

\subsection{Generative、Clustering 与 Speech Models 的共同结构}

课堂把 Jukebox、VQ-VAE、k-means、wav2vec-like codes、Tacotron 与 generative speech models 放在同一张比较图上。名称虽多，问题可以压成三列：谁产生 latent units、谁对 units 建模、谁负责重建或读出。这样才能看见“Transformer 预测 code”与“Transformer 直接产生 waveform”处在不同抽象层。

\lecturefigure{slide-33-audio-methods-comparison.jpg}{Audio methods comparison：不同系统在 encoder、discretizer、sequence model 与 decoder 上分工。}{Stanford Online 官方视频 00:31:58--00:33:21。}

\begin{table}[H]
\centering
\small
\caption{音频自监督/生成路线的机制比较。}
\begin{tabularx}{\textwidth}{L{0.20\textwidth}L{0.22\textwidth}L{0.24\textwidth}X}
\toprule
路线 & 学习信号 & 主要保留内容 & 主要风险 \\
\midrule
Contrastive & 区分 positive 与 negatives & augmentation 后仍稳定的属性 & False negatives、错误 invariance、batch dependence。 \\
Next-code prediction & 预测未来离散 codes & 对未来有预测力的状态 & Tokenizer bottleneck、exposure bias、局部捷径。 \\
Waveform prediction & 预测下一个 sample & 采样级细节与局部动力学 & 序列极长，metric 与听感错位。 \\
Reconstruction & 从 latent 恢复输入 & Decoder 能重建的内容 & 容易保留低级细节而忽略任务语义。 \\
\bottomrule
\end{tabularx}
\end{table}

\lecturefigure{slide-34-textless-nlp.jpg}{Textless NLP / generative speech：用离散 speech units 进行序列建模再解码。}{Stanford Online 官方视频 00:33:22--00:33:41。}

\subsection{预测未来为什么可能产生 Summary}

如果当前 latent state 能准确预测未来 codes，它必须压缩历史中与后续有关的信息。讲者把这一点称为“better prediction, better summarization”。形式上，希望 representation $h_t$ 使
\[
p(z_{t+1:t+K}\mid z_{\le t})
\approx p(z_{t+1:t+K}\mid h_t),
\]
其中 $z_t$ 是离散 code。这个近似表达 predictive state 的目标；它并不保证 $h_t$ 对所有未来任务都充分，也可能只学到数据中的短期可预测模式。

\lecturefigure{slide-35-long-dependency-cluster-model.jpg}{用 cluster-code sequence 学 long-term dependency：预测能力作为 representation quality 的代理。}{Stanford Online 官方视频 00:33:42--00:34:51。}

\teachervoice{“预测得越好，summary 越好”是课程的关键假设，而不是定义。若未来主要由局部纹理决定，模型可能在不理解场景语义的情况下获得好 prediction loss。}

\subsection{怎样公平比较 Learned Representations}

论文尽量固定 dataset 和 encoder family，冻结表示后训练 linear head。设表示为 $h=f_\theta(x)$，linear probe 为
\[
\hat{y}=\operatorname{softmax}(Wh+b),
\]
训练时只更新 $W,b$。若 accuracy 高，说明标签信息能被线性读取；若低，可能是信息缺失，也可能只是需要 nonlinear readout。固定 encoder/data 能减少 confounding，但 pretraining objective、augmentation、hyperparameter search 与 code rate 仍会影响结果。

\lecturefigure{slide-36-unsupervised-representation-evidence.jpg}{Controlled comparison 的设计：same encoder、same dataset、linear readout。}{Stanford Online 官方视频 00:34:52--00:36:07。}

\lecturefigure{slide-37-same-encoder-same-data-surface.jpg}{表示结果：supervised 71.3\%，contrastive 55.6\%，next-code Transformer 60.1\%。}{Stanford Online 官方视频 00:36:08--00:37:09；arXiv:2010.11459。}

\begin{knowledgebox}{读图：先看 control，再看 gap}
同 encoder、同 dataset 使 objective 成为主要变化源。Next-code Transformer 高于该 contrastive setup，但仍落后 supervised baseline 约 11.2 points。右侧训练数据/hidden-size surface 说明规模会改变误差，但不能直接证明“更多无标签数据一定消除 gap”；那是讲者提出的 scaling hypothesis。
\end{knowledgebox}

\begin{warningbox}{Linear probe 的结论边界}
71.3\%、55.6\%、60.1\% 是当前 dataset/encoder/probe 的 accuracy。它们不比较 generation quality、few-shot transfer、robustness、calibration 或 compute efficiency，也不能宣称 supervised 与 unsupervised representation 只差一个固定百分比。
\end{warningbox}

\teachervoice{Verma 明确把“更多数据会缩小 10\% gap”放在结果之后作为推测。讲义保留这个顺序：先报告受控表格，再讨论 scaling，而不是把推测写成论文已验证结论。}

\subsection{本章小结}

Contrastive learning 通过 augmentation 定义 invariance，predictive learning 通过 future codes 定义充分状态。VQ tokenizer、objective 与 probe 一起决定结果；同数据同 encoder 的 control 很重要，但仍不能把 linear accuracy 外推成通用 representation quality。

